\section{Introduction}
\label{sec:intro}

Our rig has two robot arms that fold a pair of shorts on a table. Five cameras
watch the task. One looks down on the table. The other four are tactile
sensors, one on each fingertip: each films a soft gel from behind, so touching
something shows up as a change in the gel's shading. We recorded 65
demonstrations by teleoperation and trained six models on them
(Section~\ref{sec:models}). Four are \emph{policies}, which map what the robot
sees to a plan of future arm commands. Two are \emph{world models}, which also
predict what the cameras will see next.

This report answers three questions, in order of importance.
\begin{enumerate}
  \item \textbf{How do the models differ on demonstrations they never saw?}
  This is the main comparison (Section~\ref{sec:compare}).
  \item \textbf{Does cropping the tactile images help?} Our earlier Grad-CAM
  pictures (defined in Section~\ref{sec:measures}) showed one policy looking at
  the edges of the tactile images, even before the fingers touched anything. A
  likely cause is light leaking in at the gel's edge. We tested the obvious fix,
  cropping the edges away (Section~\ref{sec:crop}). This question is secondary.
  \item \textbf{How well do the world models predict the cameras?} We score how
  well they reproduce frames they have seen and how well they predict frames
  they have not (Section~\ref{sec:wm}).
\end{enumerate}

The main findings are these.
\begin{itemize}
  \item \textbf{Both world models plan the next third of a second better than
  every policy} on unseen demonstrations.
  \item \textbf{Every model is much worse on unseen demonstrations.} The gap
  is largest for ACT, which almost memorises its training demonstrations. It is
  not caused by the session drifting over time.
  \item \textbf{Diffusion is the most accurate policy for the first steps of a
  plan, but its error grows fastest;} ACT's grows slowly.
  \item \textbf{Cropping the tactile images changes almost nothing} on any
  model, and the policies do not over-weight the image edges in the first place.
  \item \textbf{Touch matters in short bursts.} ACT nearly ignores the tactile
  cameras most of the time and leans on them heavily around contact.
  \item \textbf{A frame-averaged image score can miss a correct prediction} of
  a small moving part, such as a gripper.
\end{itemize}

Section~\ref{sec:methods} describes the models, the data and how we measure.
Section~\ref{sec:experiments} reports the experiments, and
Section~\ref{sec:discussion} discusses what we learned and what to do next.
Appendix~\ref{sec:provenance} explains how every figure and table was made.
